\section{Conclusions}
\label{sec:conclusions_dynamic}

This study resolves a question that frozen-time localization analyses cannot
answer: which admissible crystal perturbation wins after the complete
thermomechanical history has acted. Five conclusions delimit the resulting
advance.

\begin{enumerate}
\item Exact elimination of the algebraic micromorphic fields produces a
69-state finite generator without adding fictitious inertia. Generalized
Schur verification and stationary/commuting-normal positive controls show
that the observed departure from accumulated frozen growth is a physical
consequence of non-normal, rotating and noncommuting HCP tangents, rather than
a state-space, projection or implementation mismatch.

\item The answer depends on the input--output question. Across 70
selector--norm--horizon re-optimizations, unrestricted full-state and
mechanical maps can select a compact-domain high-frequency wave branch that
has no finite localization wavelength. Constitutive-response selectors
instead retain interior maxima; the baseline terminal
constitutive/constitutive map gives
$k^*=7.724\times10^3~\mathrm{m^{-1}}$ and the 1025-state map gives
$\mathcal G=6.338\times10^5$.  Its sampled near-optimal set is narrow,
whereas the onset set spans almost the complete challenged
direction--wavenumber domain and therefore does not define a unique onset
scale. Same-selector piecewise-frozen products
converge to the time-ordered propagator, and no non-rank-one case triggers the
formal 5\% singular-value degeneracy gate.

\item The terminal constitutive basin transfers across three added
rate/temperature paths, but both selected wavenumber and amplification change.
Local sensitivity further separates amplification controls from scale
controls: $\rho_{m0}$ controls gain at both horizons, $K_\Lambda$ becomes
dominant at the terminal horizon, and $\ell_\chi$ can shift onset $k^*$ with
negligible change in gain. The
Grade-II titanium comparison retains only stress-before-temperature ordering
and rejects quantitative agreement in peak time, temperature rise and
post-localization collapse. The present card is therefore a verification
seed, not a calibrated titanium predictor.

\item Mechanism attribution is time dependent: the leading cross-coupling
changes from dislocation--plastic near onset to inertia--plastic at the
terminal horizon, while full-rank slip-gradient recovery controls the sampled
quasi-static high-$k$ branch more completely than the rank-deficient Nye term.
A four-mode Galerkin replay transports both re-optimized singular inputs
within its small-amplitude acceptance gates.  The independent dealiased
16/32/64-cell ladder converges in gain, and separate time-step and
integration-format controls pass.  The raw all-mode failure is localized to
the even-grid Nyquist coordinate, but the primary terminal audit places
84.58\% of the nonlinear right-hand-side energy above the retained cutoff.  The evidence therefore
validates finite-time linear selection and resolved finite-band transport, not
complete nonlinear harmonic closure or finite-amplitude saturation.

\item The closed result is therefore finite-time, selector-dependent linear
selection on an evolving HCP path. The same-window audit still rejects a
seed-independent mature width, and the present theory does not supply
specimen calibration, propagation resistance, fracture-like dissipation or a
strict HCP-to-Bai analytical degeneration.  The open nonlinear spectral-flux
gate further shows that finite-band linear selection alone does not regularize
nonlinear harmonic transfer. Complete generators, singular
subspaces, optimization traces, parameter provenance and runtime contracts
make those open gates explicit and reproducible.
\end{enumerate}

The central distinction is consequently sharper than a new instability
number: a finite-time propagator can identify which crystal perturbation is
amplified by an evolving path, while the present evidence also states when
that perturbation cannot be interpreted as a material shear-band width.
